\begin{figure}[ht]
\centering
\begin{tikzpicture}[>=Stealth, font=\small, node distance=1.9cm]
  \node[draw, rounded corners, fill=blue!8, minimum width=4.6cm, minimum height=1cm] (local)
    {局部数据：$a_p(E)=p+1-\#E(\F_p)$};
  \node[draw, rounded corners, fill=green!8, below=of local, minimum width=4.8cm, minimum height=1cm] (L)
    {$L(s,E)$ 与中心值 $s=1$};
  \node[draw, rounded corners, fill=orange!10, below=of L, minimum width=4.8cm, minimum height=1cm] (rank)
    {全局数据：$E(\Q)\cong \Z^r\oplus E(\Q)_{\mathrm{tors}}$};
  \draw[->, thick] (local) -- node[right] {Euler 乘积} (L);
  \draw[->, thick] (L) -- node[right] {BSD 预言} (rank);
  \node[left=2.6cm of L] {$\ord_{s=1}L(s,E)$};
  \draw[dashed, ->, thick] ($(L.west)+(-0.2,0)$) -- ($(rank.west)+(-0.2,0)$);
\end{tikzpicture}
\caption{Birch--Swinnerton-Dyer 猜想把局部点数拼成的 $L$-函数，与椭圆曲线有理点群的秩联系起来。Eichler--Shimura 理论负责把左边的解析对象变得可计算。}
\label{fig:ch08-bsd-bridge}
\end{figure}
